\documentclass[reqno]{amsart}
\usepackage[T1]{fontenc}
\usepackage{amsmath,amssymb,amsthm}
\usepackage{xr-hyper}
\usepackage[hidelinks]{hyperref}
\externaldocument[dt-]{20261125T002000--definite-theta-lifts__workbook}
\externaldocument[tn-]{20261125T002100--definite-theta-lifts-and-norms__workbook}
\externaldocument[sw-]{20261125T015500--siegel-weil-in-the-smallest-case__workbook}
\externaldocument[dch-]{20261126T030700--the-determinant-of-a-chain__workbook}
\externaldocument[c2-]{20261126T050500--equal-multiplicities__workbook}
\externaldocument[c3-]{20261126T050600--the-exchange-identity__workbook}
\externaldocument[ob-]{20261126T040100--spaces-forms-and-twists__workbook}

% Chapter number; \ref{chap} gives it.
\newcounter{chap}
% Chapter numbers, generated from 20261125T012000--grh-workbook__guide.org
% by scripts/make-chapters.py.  Do not edit: change the order of the
% index and rerun the script (./preview-site.sh and CI do this).
%
% Usage in a sheet, after \newcounter{chap}:
%   % Chapter numbers, generated from 20261125T012000--grh-workbook__guide.org
% by scripts/make-chapters.py.  Do not edit: change the order of the
% index and rerun the script (./preview-site.sh and CI do this).
%
% Usage in a sheet, after \newcounter{chap}:
%   % Chapter numbers, generated from 20261125T012000--grh-workbook__guide.org
% by scripts/make-chapters.py.  Do not edit: change the order of the
% index and rerun the script (./preview-site.sh and CI do this).
%
% Usage in a sheet, after \newcounter{chap}:
%   \input{chapters}\chapnumber{<sheet ID>}
\makeatletter
\providecommand\chapters@set[2]{\@namedef{chapters@#1}{#2}}
\chapters@set{20261125T001000}{1}% zeros-to-ramp-sums
\chapters@set{20261125T001100}{2}% inert-primes-and-estimate-b
\chapters@set{20261125T005000}{3}% galois-fourier-and-averaging
\chapters@set{20261125T005500}{4}% columns-from-a-sphere
\chapters@set{20261125T000000}{5}% slope-method
\chapters@set{20261126T030100}{6}% a-twist-for-each-prime
\chapters@set{20261125T003000}{7}% the-chain-and-the-square-root
\chapters@set{20261125T010000}{8}% rank-lemma-and-zero-estimate
\chapters@set{20261125T010500}{9}% poincare-series-and-near-orthogonality
\chapters@set{20261124T223000}{10}% elliptic-curves-for-the-proof
\chapters@set{20261125T009500}{11}% fock-model-and-theta-intertwiners
\chapters@set{20261124T223100}{12}% complex-multiplication-and-inert-reduction
\chapters@set{20261125T008000}{13}% one-cm-determinant
\chapters@set{20261124T214500}{14}% supersingular-theta-divisibility
\chapters@set{20261124T214600}{15}% the-isotropic-plane
\chapters@set{20261125T013700}{16}% pairings-on-a-formal-group-of-height-two
\chapters@set{20261125T013500}{17}% hecke-nondegeneracy-on-the-tree
\chapters@set{20261126T030200}{18}% the-bound-is-attained
\chapters@set{20261125T009000}{19}% kronecker-limit-and-chowla-selberg
\chapters@set{20261125T002000}{20}% definite-theta-lifts
\chapters@set{20261125T015500}{21}% siegel-weil-in-the-smallest-case
\chapters@set{20261125T002100}{22}% definite-theta-lifts-and-norms
\chapters@set{20261125T016500}{23}% hecke-operators-and-theta-lifts
\chapters@set{20261125T016700}{24}% divisibility-along-the-chain
\chapters@set{20261125T004000}{25}% denominators-and-sha
\chapters@set{20261125T004100}{26}% tate-twists-and-sha
\chapters@set{20261126T020430}{27}% the-minus-part-by-gauss-sums
\chapters@set{20261125T014100}{28}% congruence-modules
\chapters@set{20261125T014500}{29}% eisenstein-congruences-and-selmer-classes
\chapters@set{20261125T014700}{30}% pseudocharacters-and-the-length-bound
\chapters@set{20261125T011000}{31}% galois-representations-black-box
\chapters@set{20261126T030300}{32}% the-data
\chapters@set{20261126T020210}{33}% bessel-poincare-series-and-their-gram-operators
\chapters@set{20261126T030400}{34}% the-seesaw-identity
\chapters@set{20261126T030500}{35}% divisibility-of-the-composite-lift
\chapters@set{20261126T030600}{36}% theorem-a-short-forms
\chapters@set{20261126T030700}{37}% the-determinant-of-a-chain
\chapters@set{20261126T030800}{38}% the-cost-of-one-lift
\chapters@set{20261126T030900}{39}% denominators-of-projections-to-lifts
\chapters@set{20261126T031000}{40}% twists-with-few-congruences
\chapters@set{20261126T031100}{41}% theorem-c-and-the-main-inequality
\providecommand\chapnumber[1]{%
  \@ifundefined{chapters@#1}%
    {\PackageError{chapters}{Sheet #1 is not listed in the index}%
      {List it in 20261125T012000--grh-workbook__guide.org\MessageBreak
       and run scripts/make-chapters.py.}}%
    {\edef\thechap{\@nameuse{chapters@#1}}}}
\makeatother
\chapnumber{<sheet ID>}
\makeatletter
\providecommand\chapters@set[2]{\@namedef{chapters@#1}{#2}}
\chapters@set{20261125T001000}{1}% zeros-to-ramp-sums
\chapters@set{20261125T001100}{2}% inert-primes-and-estimate-b
\chapters@set{20261125T005000}{3}% galois-fourier-and-averaging
\chapters@set{20261125T005500}{4}% columns-from-a-sphere
\chapters@set{20261125T000000}{5}% slope-method
\chapters@set{20261126T030100}{6}% a-twist-for-each-prime
\chapters@set{20261125T003000}{7}% the-chain-and-the-square-root
\chapters@set{20261125T010000}{8}% rank-lemma-and-zero-estimate
\chapters@set{20261125T010500}{9}% poincare-series-and-near-orthogonality
\chapters@set{20261124T223000}{10}% elliptic-curves-for-the-proof
\chapters@set{20261125T009500}{11}% fock-model-and-theta-intertwiners
\chapters@set{20261124T223100}{12}% complex-multiplication-and-inert-reduction
\chapters@set{20261125T008000}{13}% one-cm-determinant
\chapters@set{20261124T214500}{14}% supersingular-theta-divisibility
\chapters@set{20261124T214600}{15}% the-isotropic-plane
\chapters@set{20261125T013700}{16}% pairings-on-a-formal-group-of-height-two
\chapters@set{20261125T013500}{17}% hecke-nondegeneracy-on-the-tree
\chapters@set{20261126T030200}{18}% the-bound-is-attained
\chapters@set{20261125T009000}{19}% kronecker-limit-and-chowla-selberg
\chapters@set{20261125T002000}{20}% definite-theta-lifts
\chapters@set{20261125T015500}{21}% siegel-weil-in-the-smallest-case
\chapters@set{20261125T002100}{22}% definite-theta-lifts-and-norms
\chapters@set{20261125T016500}{23}% hecke-operators-and-theta-lifts
\chapters@set{20261125T016700}{24}% divisibility-along-the-chain
\chapters@set{20261125T004000}{25}% denominators-and-sha
\chapters@set{20261125T004100}{26}% tate-twists-and-sha
\chapters@set{20261126T020430}{27}% the-minus-part-by-gauss-sums
\chapters@set{20261125T014100}{28}% congruence-modules
\chapters@set{20261125T014500}{29}% eisenstein-congruences-and-selmer-classes
\chapters@set{20261125T014700}{30}% pseudocharacters-and-the-length-bound
\chapters@set{20261125T011000}{31}% galois-representations-black-box
\chapters@set{20261126T030300}{32}% the-data
\chapters@set{20261126T020210}{33}% bessel-poincare-series-and-their-gram-operators
\chapters@set{20261126T030400}{34}% the-seesaw-identity
\chapters@set{20261126T030500}{35}% divisibility-of-the-composite-lift
\chapters@set{20261126T030600}{36}% theorem-a-short-forms
\chapters@set{20261126T030700}{37}% the-determinant-of-a-chain
\chapters@set{20261126T030800}{38}% the-cost-of-one-lift
\chapters@set{20261126T030900}{39}% denominators-of-projections-to-lifts
\chapters@set{20261126T031000}{40}% twists-with-few-congruences
\chapters@set{20261126T031100}{41}% theorem-c-and-the-main-inequality
\providecommand\chapnumber[1]{%
  \@ifundefined{chapters@#1}%
    {\PackageError{chapters}{Sheet #1 is not listed in the index}%
      {List it in 20261125T012000--grh-workbook__guide.org\MessageBreak
       and run scripts/make-chapters.py.}}%
    {\edef\thechap{\@nameuse{chapters@#1}}}}
\makeatother
\chapnumber{<sheet ID>}
\makeatletter
\providecommand\chapters@set[2]{\@namedef{chapters@#1}{#2}}
\chapters@set{20261125T001000}{1}% zeros-to-ramp-sums
\chapters@set{20261125T001100}{2}% inert-primes-and-estimate-b
\chapters@set{20261125T005000}{3}% galois-fourier-and-averaging
\chapters@set{20261125T005500}{4}% columns-from-a-sphere
\chapters@set{20261125T000000}{5}% slope-method
\chapters@set{20261126T030100}{6}% a-twist-for-each-prime
\chapters@set{20261125T003000}{7}% the-chain-and-the-square-root
\chapters@set{20261125T010000}{8}% rank-lemma-and-zero-estimate
\chapters@set{20261125T010500}{9}% poincare-series-and-near-orthogonality
\chapters@set{20261124T223000}{10}% elliptic-curves-for-the-proof
\chapters@set{20261125T009500}{11}% fock-model-and-theta-intertwiners
\chapters@set{20261124T223100}{12}% complex-multiplication-and-inert-reduction
\chapters@set{20261125T008000}{13}% one-cm-determinant
\chapters@set{20261124T214500}{14}% supersingular-theta-divisibility
\chapters@set{20261124T214600}{15}% the-isotropic-plane
\chapters@set{20261125T013700}{16}% pairings-on-a-formal-group-of-height-two
\chapters@set{20261125T013500}{17}% hecke-nondegeneracy-on-the-tree
\chapters@set{20261126T030200}{18}% the-bound-is-attained
\chapters@set{20261125T009000}{19}% kronecker-limit-and-chowla-selberg
\chapters@set{20261125T002000}{20}% definite-theta-lifts
\chapters@set{20261125T015500}{21}% siegel-weil-in-the-smallest-case
\chapters@set{20261125T002100}{22}% definite-theta-lifts-and-norms
\chapters@set{20261125T016500}{23}% hecke-operators-and-theta-lifts
\chapters@set{20261125T016700}{24}% divisibility-along-the-chain
\chapters@set{20261125T004000}{25}% denominators-and-sha
\chapters@set{20261125T004100}{26}% tate-twists-and-sha
\chapters@set{20261126T020430}{27}% the-minus-part-by-gauss-sums
\chapters@set{20261125T014100}{28}% congruence-modules
\chapters@set{20261125T014500}{29}% eisenstein-congruences-and-selmer-classes
\chapters@set{20261125T014700}{30}% pseudocharacters-and-the-length-bound
\chapters@set{20261125T011000}{31}% galois-representations-black-box
\chapters@set{20261126T030300}{32}% the-data
\chapters@set{20261126T020210}{33}% bessel-poincare-series-and-their-gram-operators
\chapters@set{20261126T030400}{34}% the-seesaw-identity
\chapters@set{20261126T030500}{35}% divisibility-of-the-composite-lift
\chapters@set{20261126T030600}{36}% theorem-a-short-forms
\chapters@set{20261126T030700}{37}% the-determinant-of-a-chain
\chapters@set{20261126T030800}{38}% the-cost-of-one-lift
\chapters@set{20261126T030900}{39}% denominators-of-projections-to-lifts
\chapters@set{20261126T031000}{40}% twists-with-few-congruences
\chapters@set{20261126T031100}{41}% theorem-c-and-the-main-inequality
\providecommand\chapnumber[1]{%
  \@ifundefined{chapters@#1}%
    {\PackageError{chapters}{Sheet #1 is not listed in the index}%
      {List it in 20261125T012000--grh-workbook__guide.org\MessageBreak
       and run scripts/make-chapters.py.}}%
    {\edef\thechap{\@nameuse{chapters@#1}}}}
\makeatother
\chapnumber{20261126T050400}
\renewcommand{\thesection}{\thechap.\arabic{section}}

\theoremstyle{definition}
\newtheorem{problem}{Problem}[section]
\newtheorem{definition}{Definition}[section]
\theoremstyle{plain}
\newtheorem{theorem}{Theorem}[section]

\newcommand{\Z}{\mathbb{Z}}
\newcommand{\C}{\mathbb{C}}
\newcommand{\F}{\mathbb{F}}
\newcommand{\ii}{\mathrm{i}}
\newcommand{\Oc}{\mathcal{O}}
\newcommand{\Ac}{\mathcal{A}}
\newcommand{\Mch}{\mathcal{M}}
\newcommand{\Rc}{\mathcal{R}}
\newcommand{\Cl}{\mathrm{Cl}}
\newcommand{\Nm}{\operatorname{Nm}}
\newcommand{\GL}{\mathrm{GL}}
\newcommand{\U}{\mathrm{U}}
\newcommand{\Hom}{\operatorname{Hom}}
\newcommand{\vol}{\operatorname{vol}}
\newcommand{\cch}{c_{\mathrm{ch}}}
\newcommand{\uh}{\hat u}

\setlength{\emergencystretch}{3em}

\title{The norm of one lift}
\author{}
\date{November 26, 2026}

\begin{document}
\maketitle
\refstepcounter{chap}\label{chap}

\paragraph{Prerequisites.}
Adjoints (inner products are conjugate-linear in the first argument).  From Chapter~\ref{dt-chap}: the class sets \(\Cl_n\), the groups \(\Gamma_L\) and Problems~\ref{dt-pr:aut} and~\ref{dt-pr:class-one}.  From Chapter~\ref{tn-chap}: the theta function \(\Theta_{L,L'}\), the lift~\eqref{tn-eq:lift} and Problems~\ref{tn-pr:adjoint}, \ref{tn-pr:compose} and~\ref{tn-pr:rallis}; the Gaussian constant \(\vartheta_0 := \sum_{\alpha \in \Z[\ii]} e^{-\pi\lvert\alpha\rvert^2}\) of that chapter.  From Chapter~\ref{sw-chap}: Section~\ref{sw-sec:gram}.  From Chapter~\ref{ob-chap}: \(\Mch\), \(\bar G\) (Section~\ref{ob-sec:chosen}), \(W_j\), \(V_j\) (Section~\ref{ob-sec:spaces}), \(H(V)\) (Section~\ref{ob-sec:forms}), with the Petersson product (Definition~\ref{ob-def:petersson}), tame characters \(\xi\), \(Q_\xi\), \((\psi_w)\) (Section~\ref{ob-sec:twists}), \(\mu_j^0\) (Section~\ref{ob-sec:characters}) and \(\uh_{\xi,j}\) on \(H(W_j)\) (Section~\ref{ob-sec:chain}).

\medskip
Theorem~\ref{dch-thm:main} of Chapter~\ref{dch-chap} expresses the volume of the chain, relative to the trivial twist, through the corresponding ratios for \(m\) single lifts on good forms, a power of the norm of the twist and ratios of \(L\)-values.  For each single lift the proof needs \(\uh^*\uh = (\text{scalar}) \cdot (\text{Hecke series})\).  How is the norm of a theta lift computed, and where does the twist enter?

\section{The squared norm of a lift is a theta sum on the doubled space}\label{sec:doubled}

We use the definite model of Chapter~\ref{tn-chap} in weight \(0\), from \(\U(1)\) to \(\U(3)\).  By Problem~\ref{dt-pr:class-one} the class sets of \(I_1\) and \(I_3\) have one element each; let \(1\) be the constant form on \(\Cl_1\), so \(\lVert 1\rVert^2 = 1/4\).  Write \(\Theta_{(0)} \colon \Ac_{(0)}(1) \to \Ac_{(0)}(3)\) for the lift~\eqref{tn-eq:lift} in weight \(0\), and \(\Theta'_{(0)}\) for the lift in the other direction.  In weight \(0\) the kernel of~\eqref{tn-eq:lift} at \((L, L')\) is the number \(\Theta_{L,L'}(0) = \sum_{u \in \Hom_\Oc(L, L')} e^{-\pi h(u,u)}\).

\begin{problem}\label{pr:theta-I3}
\((\Theta_{(0)} 1)(I_3) = \vartheta_0^3/4 = 0.4111137902\ldots\), where \(\vartheta_0 = 1.1803405990\ldots\) is the square of \(\sum_{n \in \Z} e^{-\pi n^2}\).
\end{problem}

\begin{problem}\label{pr:norm-I3}
\[
  \frac{\lVert \Theta_{(0)} 1 \rVert^2}{\lVert 1 \rVert^2} = \frac{\vartheta_0^6}{1536} = \frac{1}{384}\cdot\frac14 \sum_{(\alpha, \beta) \in \Oc^3 \times \Oc^3} e^{-\pi(\lvert \alpha\rvert^2 + \lvert \beta\rvert^2)} = 0.0017605682\ldots ,
\]
with \(\Oc = \Z[\ii]\) and \(\lvert \alpha\rvert^2 = \sum_k \lvert \alpha_k\rvert^2\).
\end{problem}

The middle expression is a theta sum over \(\Hom_\Oc(I_1 \oplus I_1, I_3) = \Oc^3 \times \Oc^3\); in general:

\begin{problem}\label{pr:double}
For unimodular Hermitian lattices \(L, L'\),
\[
  \Theta_{L,L'}(0)\, \overline{\Theta_{L,L'}(0)} = \Theta_{L \oplus L, L'}(0) .
\]
Hence, for the weight-\(0\) lift \(\Theta_{(0)}\) from \(\Cl_n\) to \(\Cl_{n'}\) and its reverse \(\Theta'_{(0)}\), the kernel of \(\Theta'_{(0)}\Theta_{(0)}\) at \((L, L)\), \(L \in \Cl_n\), is the sum of \(\lvert\Gamma_{L'}\rvert^{-1}\, \Theta_{L \oplus L, L'}(0)\) over \(L' \in \Cl_{n'}\).
\end{problem}

The class sum models the adelic doubling integral: for a source \(V^\flat\) and a target \(V^\sharp\), a theta kernel times the conjugate of another is a theta function on \(\Hom(V^\flat \oplus (V^\flat)^-, V^\sharp)\), and the doubled source (\((V^\flat)^-\) has the negated form) is split.  Let \(V^0\) and \(V^\sharp\) be framed Hermitian spaces, in the sense of Section~\ref{ob-sec:spaces}, over a CM field with a chosen CM type and Weil data, \(V^0\) with a complete splitting into \(p\) hyperbolic planes (in this sheet \(p\) counts planes) and \(V^\sharp\) of rank \(n^\sharp > 2p\), and let \(\omega\) be the split Weil action on the test space \(\mathcal S\) of \(\Hom(V^0, V^\sharp)\).  For a rational Lagrangian \(\mathcal L\) of \(V^0\) let \(D_{\mathcal L}(\phi)\) be the value at \(0\) of \(\phi \in \mathcal S\) in the Schr\"odinger model of a splitting whose isotropic space is \(\mathcal L\).  Then the operator norms of the \(D_{\mathcal L}\) are summable, \(g^\sharp \mapsto \theta(\omega(1, g^\sharp) \phi)\) is integrable for every \(\phi\), and
\[
  \int_{[\U(V^\sharp)]} \theta(\omega(1, g^\sharp) \phi)\, dg^\sharp = \vol[\U(V^\sharp)] \sum_{\mathcal L} D_{\mathcal L}(\phi) ,
\]
where \(\vol[\U(V^\sharp)]\) is the total mass of the quotient measure.  This is the Siegel--Weil formula in Weil's range \cite[Thm.~5]{We65}.  \emph{Status}: classical; the proof establishes it in this form.  For \(V^0 = V^\flat \oplus (V^\flat)^-\) the condition is \(2n^\flat < n^\sharp\); for \(n^\flat = 1\) it first holds for \(\U(3)\), not for the \(\U(2)\) of Problem~\ref{tn-pr:rallis}.

For the summability the proof assigns to each \(\mathcal L\) a vector \(\mathbf e \in \mathbb N^p\) of dyadic exponents (with \(0 \in \mathbb N\)); with \(\lvert\mathbf e\rvert := e_1 + \dots + e_p\) and \(f\) the number of infinite places, at most \(C\, 2^{2fp\lvert\mathbf e\rvert}\) Lagrangians have a given \(\mathbf e\), and \(\lVert D_{\mathcal L}\rVert \le C' 2^{-f n^\sharp \lvert\mathbf e\rvert}\).

\begin{problem}\label{pr:majorant}
For integers \(p, f \ge 1\) and \(n^\sharp\), \(\sum_{\mathbf e \in \mathbb N^p} 2^{-f(n^\sharp - 2p)\lvert\mathbf e\rvert} = \sum_{l \ge 0} \binom{l + p - 1}{p - 1} 2^{-f(n^\sharp - 2p) l}\).  It converges if and only if \(n^\sharp > 2p\), with sum \((1 - 2^{-f(n^\sharp - 2p)})^{-p}\); for \(p = f = 1\), \(n^\sharp = 3\) the sum is \(2\), and the terms with \(l \le 5\) give \(63/32\).
\end{problem}

\section{Only graphs survive, and unfolding}\label{sec:unfold}

Which Lagrangians of the doubled space contribute?  Over a finite field, let \(\F_9 = \F_3[\ii]\), with conjugation \(z \mapsto \bar z := z^3\).

\begin{problem}\label{pr:F9}
The norm \(z \mapsto z\bar z\) from \(\F_9\) to \(\F_3\) has fibres of sizes \(1, 4, 4\) over \(0, 1, 2\).  The form \(z_1\bar z_1 - z_2\bar z_2\) on \(\F_9^2\) has \(33\) isotropic vectors (including \(0\)) and \(4\) isotropic lines, and these lines are the graphs \(\{(\delta z, z) : z \in \F_9\}\) of the \(4\) elements \(\delta\) of norm \(1\).
\end{problem}

A \emph{frame} of a Lagrangian is a matrix whose columns are a basis of that Lagrangian.

\begin{problem}\label{pr:graph}
Let \(V^\flat\) be a nondegenerate Hermitian space of dimension \(n^\flat\) over a field with involution, with Gram matrix \(J\), and give \(V^\flat \oplus (V^\flat)^-\) the Gram matrix \(\operatorname{diag}(J, -J)\).  Let \(\mathcal L\) be a Lagrangian, with frame \(\binom{K_1}{K_2}\): a \(2n^\flat \times n^\flat\) matrix of rank \(n^\flat\) whose columns span \(\mathcal L\).  If \(K_2\) is invertible, then \(\delta := -K_1K_2^{-1}\) is an isometry of \(V^\flat\) and \(\mathcal L = \{(\delta z, -z) : z \in V^\flat\}\).  If \(V^\flat\) has no nonzero isotropic vector, every Lagrangian is of this form.
\end{problem}

For the clean data \(\phi\) of Theorem~\ref{thm:rallis} below, the proof shows that only graphs of rational isometries contribute (by Problem~\ref{pr:graph} when \(V^\flat\) is definite): for the doubled datum of the translates of \(\phi'\), \(\phi\) by \(g'\), \(g\), the term of \(\{(\delta^{-1} z, -z)\}\) is \(\langle \phi', \omega(g'^{-1}\delta g)\phi\rangle\).  Integrating against \(y(g)\overline{y'(g')}\) then unfolds, as in this finite model.

\begin{problem}\label{pr:unfold}
Let \(\Gamma\) be a subgroup of a finite group \(G\), let \(K\) be a function on \(G\) (in this sheet \(K\) is a function; the field \(\mathbb{Q}(\sqrt{-d})\) does not occur), and let \(y, y'\) be functions on \(G\) invariant on the left under \(\Gamma\), with \(\langle y', y\rangle := \sum_{g \in \Gamma\backslash G} \overline{y'(g)}\, y(g)\).  Then
\[
  \sum_{g, g' \in \Gamma\backslash G}\ \sum_{\delta \in \Gamma} K(g'^{-1}\delta g)\, y(g)\, \overline{y'(g')} = \sum_{g \in G} K(g)\, \bigl\langle y', y(\cdot\, g)\bigr\rangle .
\]
\end{problem}

The result is the Rallis inner product formula.  Let \(V^\flat\), \(V^\sharp\) be framed Hermitian spaces over a CM field with a chosen CM type and Weil data, of ranks \(n^\flat\), \(n^\sharp\) with \(2n^\flat < n^\sharp\), and let \(G^\flat\) be the adelic unitary group of \(V^\flat\).  Let \(\chi_1, \chi_2\) be continuous unitary characters of the norm-one ideles trivial on rational points, and let \(\phi, \phi'\) lie in the adelic test space of \(\Hom(V^\flat, V^\sharp)\).  Suppose that \(\phi\) is \emph{clean}: either all signs of \(V^\flat\) agree, or \(\phi = \phi_1 \otimes \phi_0\) for a Schwartz--Bruhat function \(\phi_1\) at a set of finite places disjoint from its conjugate such that, at some finite place, every matrix in the support of \(\phi_1\) has a left inverse.  Let \(y, y'\) be automorphic functions on \(G^\flat\) of rapid decay, let \(\Theta[\phi]\, y\) be the theta lift of \(y\) to the group of \(V^\sharp\) with datum \(\phi\) and characters \((\chi_1, \chi_2)\), and put \(K(g) := \langle \phi', \omega(g) \phi\rangle\, \langle y', y(\cdot\, g)\rangle\), where \(\omega(g)\) is the action of \((g, 1)\) with characters \((1, \chi_2)\) and the second factor is the Petersson product.

\begin{theorem}[Rallis inner product formula]\label{thm:rallis}
\(K\) is integrable on all of \(G^\flat\), and \(\langle \Theta[\phi'] y', \Theta[\phi] y\rangle = \vol[\U(V^\sharp)] \int_{G^\flat} K(g)\, dg\).
\end{theorem}

This is the unitary form, used here, of Rallis's inner product formula; for the classical method see \cite[\S\S\,2--3]{Ra87}.  \emph{Status}: classical; the proof derives it in this form from the Siegel--Weil formula, as Problems~\ref{pr:graph} and~\ref{pr:unfold} model.

\section{The sharp datum and the twist factor}\label{sec:sharp}

For \(1 \le j \le m\) the single lift \(\uh_{\xi,j}\) goes from \(W_j\), of rank \(r\), to \(V_j\), of rank \(r + t > 2r\); \(H(W_j)\) is finite-dimensional (Definition~\ref{ob-def:petersson}), so \(\uh_{\xi,j}\) has an adjoint.  At the places of \(\Mch \cup Q_\xi\) every residue matrix in the support of its finite datum has a left inverse; this support count and the count of all such matrices (through the target level index) give the finite factor in Theorem~\ref{thm:one-lift}.

\begin{problem}\label{pr:columns}
Of the \(5^3\) columns of height \(3\) over \(\F_5\), exactly \(124\) have a left inverse; of the \(5^6\) columns of height \(6\), exactly \(15624\).
\end{problem}

\begin{problem}\label{pr:left-inverse}
For integers \(r \ge 1\) and \(t \ge 0\), the number of \((r + t) \times r\) matrices over \(\F_q\) with a left inverse is \(q^{r(r+t)} \prod_{v = t+1}^{t+r} (1 - q^{-v})\).
\end{problem}

\begin{definition}\label{def:twist-factor}
For a finite set \(Q\) of finite places of \(E\), with residue fields \(k_w\) of order \(q_w\), define the \emph{twist factor} by
\[
  c_Q := (\Nm Q)^{rt} \prod_{w \in Q} \prod_{v = t+1}^{t+r} \bigl(1 - q_w^{-v}\bigr), \qquad \Nm Q := \prod_{w \in Q} q_w ,
\]
where \(r\) is the source rank and \(t\) the target rank minus the source rank; for the single lifts of the data, \(t = 5r\) and \(Q = Q_\xi\).
\end{definition}

\begin{problem}\label{pr:cQ}
\(c_Q = \prod_{w \in Q} (\lambda_w / q_w^{r^2})\), where \(\lambda_w\) (in this sheet) is the number of \((r+t) \times r\) matrices over \(k_w\) with a left inverse; in particular \(c_\varnothing = 1\).  Formally evaluating the expression of Definition~\ref{def:twist-factor} at \(Q = \{w\}\), \(q_w = 5\), \(r = 1\) gives \(c_Q = 124/5\) for \(t = 2\) and \(c_Q = 15624/5\) for \(t = 5\).
\end{problem}

Give the standard lattice stabilizer in each finite adelic unitary group Haar volume \(1\).  When both lifts use the same sharp datum (below) and the unprimed form has principal level, the finite integration in Theorem~\ref{thm:rallis} is a sum over cosets of the subgroup integral outside \(\Mch \cup Q_\xi\) and their conjugates, and the integral on each coset supplies the mass of the finite datum as a factor; here is the finite model.

\begin{problem}\label{pr:collapse}
Let \(\Gamma\) and \(\mathcal J\) be subgroups of a finite group \(G\), let \(y, y'\) be functions on \(\Gamma\backslash G\), and let \(y\) be invariant on the right under \(\mathcal J\).  With the pairing of Problem~\ref{pr:unfold}, \(\sum_{g \in G} \mathbf 1_{\mathcal J}(g)\, \langle y', y(\cdot\, g)\rangle = \#\mathcal J\, \langle y', y\rangle\).
\end{problem}

The lift \(\uh_{\xi,j}\) uses a highest-vector Fock datum \(\phi_\infty\) at the infinite places and a finite Weil vector \(F_{\mathrm{fin}}\), obtained from the vacuum by imposing, modulo the product of the places of \(\Mch \cup Q_\xi\), the following residue function \(F\) (in this sheet) on families \(\mathbf X = (X_w)_{w \in \Mch \cup Q_\xi}\) of \((r+t) \times r\) matrices over the \(k_w\): \(F(\mathbf X) := \prod_{w \in Q_\xi} \psi_w(\det X_w')\) if \(X_w\) is the standard inclusion \(\binom{1}{0}\) at every \(w \in \Mch\), and \(F(\mathbf X) := 0\) otherwise; here \(X_w'\) is the top \(r \times r\) block of \(X_w\), and \(\psi_w(0) := 0\).  Let \(J_{\Mch}\) (in this sheet) be the principal level of \(H(W_j)\) and \(K^\circ\) the stabilizer of the standard lattice.  Let \(S_0\) be the set of places of \(\Mch \cup Q_\xi\) and their conjugates.  In the matrix coefficients of this lift take \(\chi_2 = b_j := \mu_j^0\xi\,\upsilon^{\lfloor (r+t)/2\rfloor}\) in Theorem~\ref{thm:rallis}, with \(\upsilon\) the splitting character on the norm-one ideles.  The proof shows that \(F_{\mathrm{fin}}\) is \emph{sharp} for \(J_{\Mch}\) off \(S_0\): for every Fock vector \(v\), \(v \otimes F_{\mathrm{fin}}\) is clean and a vacuum off \(S_0\); \(\lVert F_{\mathrm{fin}}\rVert^2 = \langle F_{\mathrm{fin}}, F_{\mathrm{fin}}\rangle \ne 0\); and for finite \(g\) integral outside \(S_0\), \(\langle F_{\mathrm{fin}}, \omega(g) F_{\mathrm{fin}}\rangle = \lVert F_{\mathrm{fin}}\rVert^2\, \mathbf 1_{J_{\Mch}}(g)\).  The \emph{mass} of \(F_{\mathrm{fin}}\) is \(\lVert F_{\mathrm{fin}}\rVert^2/[K^\circ : J_{\Mch}]\), where \(\lVert F_{\mathrm{fin}}\rVert^2\) is the mean of \(\lvert F(\mathbf X)\rvert^2\) over all families \(\mathbf X\) and \([K^\circ : J_{\Mch}] = \#\bar G\).

\begin{problem}\label{pr:sum-block}
\(\sum_{\mathbf X} \lvert F(\mathbf X)\rvert^2 = \#\GL_r(Q_\xi) \cdot (\Nm Q_\xi)^{rt}\), where \(\#\GL_r(Q) := \prod_{w \in Q} \#\GL_r(k_w)\).
\end{problem}

The class-sum product on the target of \(\uh_{\xi,j}\) is the index of the level of \(V_j\) times the dimension of the coefficient representation times the Petersson product of the associated adelic forms, and the proof shows that this index times \(\#\GL_r(Q_\xi)\) is \(\prod_{w \in \Mch \cup Q_\xi} \lambda_w\), the size of the orbit of the standard inclusion.  Put
\[
  \cch := \prod_{w \in \Mch} \biggl( \frac{1}{\#\GL_r(k_w)} \prod_{v = t+1}^{t+r} \bigl(1 - q_w^{-v}\bigr) \biggr) ,
\]
a nonzero number that does not depend on \(\xi\).

\begin{problem}\label{pr:mass}
With \(\lambda_w\) as in Problem~\ref{pr:cQ},
\[
  \frac{(\Nm Q_\xi)^{rt}}{\#\bar G} \prod_{w \in \Mch \cup Q_\xi} \frac{\lambda_w}{q_w^{(r+t)r}} = \cch\, c_{Q_\xi} .
\]
So the index of the level of \(V_j\) times the mass of \(F_{\mathrm{fin}}\) is \(\cch c_{Q_\xi}\), and only the factor \(c_{Q_\xi}\) depends on \(\xi\).
\end{problem}

At the other places the integral becomes a sum over Hecke cosets.  Write \(A\) (in this sheet) for the set of finite places outside \(S_0\).  For a finite point \(\gamma\) of the unitary group of a framed space \(V\) of rank \(r\) let \(\Delta_f(\gamma)\) be its finite Weil multiplier and (in this sheet)
\[
  \mathrm{mc}(\gamma) := \bigl(b_j(\det\gamma)\, \Delta_f(\gamma)^{r+t}\bigr)^{-1} ,
\]
so \(\mathrm{mc}(1) = 1\).  For a datum that is a vacuum on \(A\), left multiplication by \(\gamma\) supported on \(A\) multiplies its matrix coefficient at \(g\), with \(g\) integral on \(A\), by \(\mathrm{mc}(\gamma)\).

\begin{definition}\label{def:rallis-operator}
Let \(j \ge 0\) be an integer, \(\xi\) a tame character, and \(V\) a framed space of rank \(r\), with any signs, such as \(W_j\) or \(Y_0\).  Define the \emph{Rallis operator} \(\Rc_j(\xi; V)\) as the endomorphism of \(H(V)\) given by
\[
  \Rc_j(\xi; V)\, y := \sum_\gamma \mathrm{mc}(\gamma)\, y(\cdot\, \gamma) ,
\]
where \(\gamma\) runs over representatives, equal to \(1\) outside \(A\), of the finite points of \(\U(V)\) modulo those integral at every place of \(A\).
\end{definition}

Through the frames, all framed spaces \(V\) of rank \(r\) have the same finite adelic unitary group, so one family of representatives \(\gamma\) serves every \(V\); the representatives are \(1\) at the places of \(\Mch\), and neither they nor the coefficients \(\mathrm{mc}(\gamma)\) depend on \(V\).

The proof shows that for \(2r < r + t\) the series converges absolutely, lies in \(H(V)\), and \(\langle y', \Rc_j(\xi; V)\, y\rangle = \sum_\gamma \mathrm{mc}(\gamma) \langle y', y(\cdot\,\gamma)\rangle\) for \(y, y' \in H(V)\); it proves the summability of the coefficients by counting neighbouring lattices.  For \(1 \le j \le m\), the sharp-data form of Theorem~\ref{thm:rallis} and Problem~\ref{pr:mass} give, for every tame character \(\xi\) and all \(y, y' \in H(W_j)\),
\[
  \langle \uh_{\xi,j}\, y', \uh_{\xi,j}\, y\rangle = a_j\, c_{Q_\xi}\, \langle y', \Rc_j(\xi; W_j)\, y\rangle ,
\]
where \(a_j\) is \(\cch\) times the dimension of the coefficient representation times a constant \(c_\infty\) (in this sheet) that depends only on the spaces and the real data.  As \(\xi\) has weight \(0\), \(a_j\) does not depend on \(\xi\).  Here \(c_\infty = \vol[\U(V_j)]\, z_\infty\), where \(z_\infty\) (in this sheet) is the integral over the real source group of \(\langle \phi_\infty, \Omega(x,1)\phi_\infty\rangle\, \overline{e(x)}/e(1)\), with \(e\) the source type kernel and \(\Omega\) the real Weil action.  Cartan averaging on the lowest compact type gives positive real part to the pairing of each local factor of \(\phi_\infty\) with its Weil average against the local discrete-series kernel.  The integral \(z_\infty\) is the product of these pairings divided by \(e(1) \ne 0\), so \(z_\infty \ne 0\).  The proof also shows that the lift with datum \(\phi_\infty \otimes F_{\mathrm{fin}}\) takes forms of the source type to forms of the target type; for the correspondence at infinity see \cite{Li90}.

\section{Injectivity and the theorem}\label{sec:theorem}

Is \(\Rc_j(\xi; V)\) injective?  The proof partitions \(A\) into finitely many disjoint conjugation-stable sets, on each of which the neighbour weights (the number of steps at \(w\) times \(q_w^{-(r+t)/2}\)) add up to at most \(1/3\).  By sharpness, \(b_j \circ \det\) is trivial on the integral points supported on \(A\).  On bounded functions fixed by the integral points supported on these places, the series over a union of parts is the composite of their series; on one part the moduli of the \(\mathrm{mc}(\gamma)\) off the trivial coset add up to at most \(1/2\), by a walk whose one-step weights total at most \(1/3\).

\begin{problem}\label{pr:half}
\(\sum_{l \ge 1} 3^{-l} = 1/2\).  Let \(y\) be a bounded function on a group \(G\), and let \(g_i \in G\) and \(\zeta_i \in \C\), \(i \ge 1\), with \(\sum_i \lvert \zeta_i\rvert \le 1/2\).  If \(y + \sum_i \zeta_i\, y(\cdot\, g_i) = 0\), then \(y = 0\).
\end{problem}

With Problem~\ref{pr:half} this gives what the proof shows: \(\Rc_j(\xi; V)\) is injective.

\begin{problem}\label{pr:injective}
Let \(u \colon \mathcal E \to \mathcal E'\) be a linear map of finite-dimensional complex inner product spaces, \(\Rc\) an injective endomorphism of \(\mathcal E\), and \(\zeta \in \C\), \(\zeta \ne 0\).  If \(\langle u y', u y\rangle = \zeta \langle y', \Rc y\rangle\) for all \(y, y' \in \mathcal E\), then \(u\) is injective.
\end{problem}

\begin{problem}\label{pr:operator}
In Problem~\ref{pr:injective}, the identity \(\langle u y', u y\rangle = \zeta \langle y', \Rc y\rangle\) for all \(y, y'\) is equivalent to \(u^* u = \zeta\, \Rc\).
\end{problem}

Problems~\ref{pr:collapse}--\ref{pr:mass} compute the finite factor \(\cch c_{Q_\xi}\) of the Gram identity of Section~\ref{sec:sharp}, whose other inputs we took from the proof, and Problem~\ref{pr:operator} turns it into the following theorem, which the proof states in this form.

\begin{theorem}\label{thm:one-lift}
For every \(1 \le j \le m\) there is \(a_j \in \C\), \(a_j \ne 0\), such that for every tame character \(\xi\),
\[
  \uh_{\xi,j}^*\, \uh_{\xi,j} = a_j\, c_{Q_\xi}\, \Rc_j(\xi; W_j)
\]
on \(H(W_j)\).
\end{theorem}

By Problem~\ref{pr:injective} and the injectivity of \(\Rc_j(\xi; W_j)\), the lift \(\uh_{\xi,j}\) is injective.

Taken as given were the Siegel--Weil formula in Weil's range, Theorem~\ref{thm:rallis} in its sharp form, the summability of the coefficients, the archimedean properties of \(\phi_\infty\), the sharpness of \(F_{\mathrm{fin}}\), the class-sum product and the index identity on the target of \(\uh_{\xi,j}\), and the change of matrix coefficients under left multiplication, and the bound on the neighbour weights behind the injectivity of \(\Rc_j(\xi; V)\).  Chapter~\ref{c2-chap} compares \(\det \Rc_j(\xi; W_j)\) and \(\det \Rc_j(\xi; Y_0)\) on good forms.  The rank hypothesis of Theorem~\ref{thm:rallis} holds for \(Y_0 \to Y_1\) but not for \(Y_1 \to Y_2\), as \(2(r + t) > r + 2t\); the chain identity uses the exchange identity of Chapter~\ref{c3-chap} instead.

\paragraph{In the Lean formalization.}
\begin{sloppypar}
In the Lean proof \cite{GL}, Theorem~\ref{thm:one-lift} is \texttt{New.Data.exists\_\allowbreak adjoint\_\allowbreak singleLift\_\allowbreak comp\_\allowbreak eq} in \texttt{New/\allowbreak Determinant/\allowbreak Norms/\allowbreak AdjointThetaLift.lean}, and the Gram identity of Section~\ref{sec:sharp}, with \(\langle y', \Rc_j(\xi; W_j)\, y\rangle\) written as its series over \(\gamma\), is \texttt{Data.\allowbreak exists\_\allowbreak hasSum\_\allowbreak inner\_\allowbreak singleLift} there.  See also, in that file, \texttt{DescentPair.\allowbreak twistFactor}, \texttt{Data.\allowbreak rallisOperator}, \texttt{DescentPair.\allowbreak isSharp\_\allowbreak blockDatum} and \texttt{DescentPair.\allowbreak gramFactor\_\allowbreak div\_\allowbreak card}, and in \texttt{Lib/\allowbreak SiegelWeil/} the declarations \texttt{New.\allowbreak SiegelWeil.\allowbreak siegel\_\allowbreak weil}, \texttt{New.\allowbreak SiegelWeil.\allowbreak inner\_\allowbreak product\_\allowbreak formula} (Theorem~\ref{thm:rallis}), \texttt{New.\allowbreak SiegelWeil.\allowbreak rallis\_\allowbreak inner\_\allowbreak product\_\allowbreak cosetTerm} and \texttt{New.\allowbreak SiegelWeil.\allowbreak eq\_\allowbreak zero\_\allowbreak of\_\allowbreak rallisSeries\_\allowbreak eq\_\allowbreak zero}.  The nonvanishing of \(z_\infty\) is \texttt{New.integral\_\allowbreak fockInner\_\allowbreak mul\_\allowbreak typeKernel\_\allowbreak ne\_\allowbreak zero} in \texttt{Lib/Weil/Lift.lean}, using \texttt{New.OscillatorPair.re\_\allowbreak inner\_\allowbreak integratedRep\_\allowbreak pos} in \texttt{Lib/DiscreteSeries/Theta/LocalCorrespondence.lean}.
\end{sloppypar}

\section*{Hints}

\subsection*{\S\ref{sec:doubled}}
\begin{description}
\item[\ref{pr:theta-I3}] Equation~\eqref{tn-eq:lift} with \(\lvert\Gamma_{I_1}\rvert = 4\); \(\Hom_\Oc(I_1, I_3) = \Oc^3\), and the sum over \(\Oc = \Z \oplus \Z\ii\) factors.
\item[\ref{pr:norm-I3}] The squared norm on \(\Ac_{(0)}(3)\) and Problem~\ref{dt-pr:aut} for \(I_3\).
\item[\ref{pr:double}] \(\Hom_\Oc(L \oplus L, L') = \Hom_\Oc(L, L')^2\); Problems~\ref{tn-pr:compose} and~\ref{tn-pr:adjoint}.
\item[\ref{pr:majorant}] Count the \(\mathbf e \in \mathbb N^p\) with \(\lvert\mathbf e\rvert = l\), and expand \((1 - z)^{-p}\).
\end{description}

\subsection*{\S\ref{sec:unfold}}
\begin{description}
\item[\ref{pr:F9}] \((u_1 + u_2\ii)\overline{(u_1 + u_2\ii)} = u_1^2 + u_2^2\).
\item[\ref{pr:graph}] Isotropy reads \(K_1^* J K_1 = K_2^* J K_2\); then look at vectors of \(\mathcal L\) of the form \((z, 0)\).
\item[\ref{pr:unfold}] Fix representatives of \(\Gamma\backslash G\) and the representative \(g'\).  As \(\delta\) and the representative \(g\) vary together, which values does \(g'^{-1}\delta g\) take?
\end{description}

\subsection*{\S\ref{sec:sharp}}
\begin{description}
\item[\ref{pr:left-inverse}] Left inverse means independent columns; choose them one at a time.
\item[\ref{pr:sum-block}] Sum place by place; \(\lvert\psi_w(d)\rvert = 1\) for \(d \ne 0\).
\item[\ref{pr:mass}] Problem~\ref{pr:left-inverse}, and the definition of \(\bar G\).
\end{description}

\subsection*{\S\ref{sec:theorem}}
\begin{description}
\item[\ref{pr:half}] Compare \(\sup \lvert y\rvert\) with its half.
\end{description}

{\small
\begin{thebibliography}{Ra87}
\bibitem[GL]{GL} L.~Sernau, \emph{GRH for Dirichlet \(L\)-functions, in Lean~4}, a shorter rewriting of the Lean proof, text by Claude, \url{https://github.com/lsernau/grh-lean}, commit \texttt{72fead1}, readable edition \texttt{grh/}.
\bibitem[Li90]{Li90} J.-S.~Li, \emph{Theta lifting for unitary representations with nonzero cohomology}, Duke Math.\ J.\ \textbf{61} (1990), 913--937.
\bibitem[Ra87]{Ra87} S.~Rallis, \emph{\(L\)-functions and the oscillator representation}, Lecture Notes in Math.\ 1245, Springer, 1987.
\bibitem[We65]{We65} A.~Weil, \emph{Sur la formule de Siegel dans la th\'eorie des groupes classiques}, Acta Math.\ \textbf{113} (1965), 1--87.
\end{thebibliography}
}

\end{document}
